% figures/fig3_rama_b.tex
\begin{tikzpicture}[
    >=Stealth,
    % Estilo idéntico al bloque principal de Figuras 1 y 2
    block/.style={
        rectangle,
        draw=black!70,
        fill=black!2,
        rounded corners=4pt,
        line width=0.8pt,
        text width=5.8cm,
        align=center,
        inner sep=8pt,
        minimum height=1.35cm,
        font=\small\sffamily
    },
    % Fórmulas matemáticas de la columna derecha
    state/.style={
        font=\small,
        align=center,
        inner sep=6pt
    },
    % Flechas principales
    arrow/.style={
        ->,
        line width=1.05pt,
        draw=black!75,
        shorten >= 2pt,
        shorten <= 2pt
    },
    % Operador circular idéntico a Figuras 1 y 2
    oplabel/.style={
        circle,
        draw=black!85,
        fill=white,
        line width=0.9pt,
        inner sep=3.5pt,
        font=\normalsize\bfseries
    }
]

    % =========================================================================
    % Nivel 0: Entrada de la Rama B (y = 0.0)
    % =========================================================================
    \node[block] (inB) at (0, 0) {
        \textbf{Vector de estado $\vec{P}$}\\
        \footnotesize Entrada a la Rama B (Espacio recíproco en $\ell^2$)
    };
    \node[state] (s0) at (7.2, 0) {
        $\vec{P} = \begin{pmatrix} \vec{p}_1 \\ \vec{p}_2 \end{pmatrix} \in \mathbb{R}^N \times \mathbb{R}^N$
    };

    % Operador 1
    \node[oplabel] (op1) at (7.2, -1.8) {$\hat{\mathcal{F}}_{\text{FFT}}$};

    % =========================================================================
    % Nivel 1: Espectro Cartesiano (y = -3.6)
    % =========================================================================
    \node[block] (b1) at (0, -3.6) {
        \textbf{Transformada discreta de Fourier}\\
        \footnotesize Vectores espectrales cartesianos $\{X_R, X_I\}$
    };
    \node[state] (s1) at (7.2, -3.6) {
        $\mathbf{Z}[k] = X_R[k] + j X_I[k] \in \mathbb{C}^{N/2}$
    };

    % Operador 2
    \node[oplabel] (op2) at (7.2, -5.4) {$\hat{\mathcal{R}}_C$};

    % =========================================================================
    % Nivel 2: Espectro Polar / CORDIC (y = -7.2)
    % =========================================================================
    \node[block] (b2) at (0, -7.2) {
        \textbf{Rotaciones polares CORDIC}\\
        \footnotesize Transformación afín de Givens
    };
    \node[state] (s2) at (7.2, -7.2) {
        $\left( \norm{\mathbf{Z}[k]}, \, \Phi[k] \right)$
    };

    % Operador 3
    \node[oplabel] (op3) at (7.2, -9.0) {$\hat{\mathcal{M}}_{\text{sinc}}$};

    % =========================================================================
    % Nivel 3: Superresolución Portadora (y = -10.8)
    % =========================================================================
    \node[block] (b3) at (0, -10.8) {
        \textbf{Estimación analítica sub-Hertz}\\
        \footnotesize Inversión no lineal del lóbulo sinc
    };
    \node[state] (s3) at (7.2, -10.8) {
        $\left( f_0 \pm \Delta f_0, \, \norm{\mathbf{Z}(f_0)}, \, \Phi(f_0) \right)$
    };

    % =========================================================================
    % CONEXIONES VERTICALES
    % =========================================================================
    \draw[arrow] (inB) -- (b1);
    \draw[arrow] (b1) -- (b2);
    \draw[arrow] (b2) -- (b3);

    \draw[arrow] (s0) -- (op1);
    \draw[arrow] (op1) -- (s1);

    \draw[arrow] (s1) -- (op2);
    \draw[arrow] (op2) -- (s2);

    \draw[arrow] (s2) -- (op3);
    \draw[arrow] (op3) -- (s3);

\end{tikzpicture}